% Generated from locale/tr/content/first-order-logic/tableaux/quantifier-rules.tex; do not hand-edit.


\olfileid[tr]{fol}{tab}{qrl}

\olsection{Niceleyici Kuralları}

\subsection{$\lforall$ Kuralları}

\begin{defish}
\AxiomC{\sFmla{\True}{\lforall[x][!A(x)]}}
\RightLabel{\TRule{\True}{\lforall}}
\UnaryInfC{\sFmla{\True}{!A(t)}}
\DisplayProof
\hfill
\AxiomC{\sFmla{\False}{\lforall[x][!A(x)]}}
\RightLabel{\TRule{\False}{\lforall}}
\UnaryInfC{\sFmla{\False}{!A(a)}}
\DisplayProof
\end{defish}

\TRule{\True}{\lforall} kuralında $t$ kapalı bir terimdir (yani değişken
içermeyen bir terimdir). \TRule{\False}{\lforall} kuralında $a$,
\TRule{\False}{\lforall} kuralının üstündeki dalın hiçbir yerinde
geçmemesi gereken !!a{constant}'tır. $a$'ya
\TRule{\False}{\lforall} çıkarımının \emph{özdeğişkeni} deriz.\footnote{Yukarıdaki
kuralda $a$ bir !!a{constant} olduğu hâlde ``özdeğişken'' terimini
kullanıyoruz. Bunun tarihsel nedenleri vardır.}

\subsection{$\lexists$ Kuralları}

\begin{defish}
\AxiomC{\sFmla{\True}{\lexists[x][!A(x)]}}
\RightLabel{\TRule{\True}{\lexists}}
\UnaryInfC{\sFmla{\True}{!A(a)}}
\DisplayProof
\hfill
\AxiomC{\sFmla{\False}{\lexists[x][!A(x)]}}
\RightLabel{\TRule{\False}{\lexists}}
\UnaryInfC{\sFmla{\False}{!A(t)}}
\DisplayProof
\end{defish}

Yine $t$ kapalı bir terimdir; $a$ ise~\TRule{\True}{\lexists} kuralının
üstündeki dalda geçmeyen !!a{constant}'tır. $a$'ya
\TRule{\True}{\lexists} çıkarımının \emph{özdeğişkeni} deriz.

Bir özdeğişkenin \TRule{\False}{\lforall} ya da
\TRule{\True}{\lexists} çıkarımının üstündeki dalda geçmemesi koşuluna
\emph{özdeğişken koşulu} denir.

\begin{explain}
$!A$ bir !!a{formula} olup !!{variable}~$x$ onda serbest geçtiğinde bunu
~$!A(x)$ yazarak belirttiğimizi anımsayın. Aynı bağlamda $!A(t)$,
~$\Subst{!A}{t}{x}$ için bir kısaltmadır. Dolayısıyla
\TRule{\False}{\lexists} kuralını şöyle de yazabilirdik:
\begin{prooftree}
  \AxiomC{\sFmla{\False}{\lexists[x][!A]}}
  \RightLabel{\TRule{\False}{\lexists}}
  \UnaryInfC{\sFmla{\False}{\Subst{!A}{t}{x}}}
\end{prooftree}
$t$'nin~$!A$ içinde zaten geçebileceğine dikkat edin; örneğin $!A$,
~$\Atom{\Obj P}{t,x}$ olabilir. Bu nedenle
~$ \sFmla{\False}{\lexists[x][\Atom{\Obj P}{t,x}]}$ ifadesinden
$\sFmla{\False}{\Atom{\Obj P}{t,t}}$ çıkarmak
~\TRule{\False}{\lexists} kuralının doğru bir uygulamasıdır. Buna karşılık
\TRule{\False}{\lforall} ve~\TRule{\True}{\lexists} kurallarındaki
özdeğişken koşulları, !!{constant}~$a$'nın $!A$ içinde geçmemesini
gerektirir. Dolayısıyla
$\sFmla{\False}{\lforall[x][\Atom{\Obj P}{a,x}]}$ ifadesinden
$\sFmla{\False}{\Atom{\Obj P}{a,a}}$ ifadesini
~$\TRule{\False}{\lforall}$ kullanarak doğru biçimde çıkaramazsınız.
\end{explain}

\begin{explain}
\TRule{\True}{\lforall} ve \TRule{\False}{\lexists} kurallarında
terim~$t$ herhangi bir kapalı terim olabilir. Buna karşılık
\TRule{\True}{\lexists} ve \TRule{\False}{\lforall} kurallarında
özdeğişken koşulu, !!{constant}~$a$'nın ilgili çıkarımın üstündeki dalların
hiçbir yerinde geçmemesini gerektirir. Bu koşul, sistemin sağlam olmasını
güvence altına almak için gereklidir. Koşul olmasaydı aşağıdaki ağaç,
$\lexists[x][\formula{A}(x)] \lif \lforall[x][\formula{A}(x)]$ için kapalı
bir !!{tableau} olurdu:
\begin{center}
\begin{tableau}{}
  [\sFmla{\False}{\lexists[x][\formula{A}(x)] \lif \lforall[x][\formula{A}(x)]}, just=\TAss
    [\sFmla{\True}{\lexists[x][\formula{A}(x)]},
      just={\TRule{\False}{\lif}[1]}
      [\sFmla{\False}{\lforall[x][\formula{A}(x)]},
        just={\TRule{\False}{\lif}[1]}
        [\sFmla{\True}{\formula{A}(a)},
          just={\TRule{\True}{\lexists}[2]}
          [\sFmla{\False}{\formula{A}(a)},
            just={\TRule{\False}{\lforall}[3]}, close]
        ]
      ]
    ]
  ]
\end{tableau}
\end{center}
Oysa $\lexists[x][\formula{A}(x)] \lif
\lforall[x][\formula{A}(x)]$ geçerli değildir.
\end{explain}

